% =====================================================================
% PARTE I — Le fondamenta
% =====================================================================
\part{Le fondamenta}

\noindent\itshape Cosa assumi prima di scrivere la prima regola.
Cinque capitoli che fissano il vocabolario: il rapporto tra formalismo e materia
(\S\ref{cap:partenza}); i quattro giudizi del metalinguaggio (\S\ref{cap:giudizi});
la grammatica di contesti e dell'operazione $\vdash$ (\S\ref{cap:contesti});
il test operativo per riconoscere linguaggio e metalinguaggio
(\S\ref{cap:test}); lo schema F-I-E-C che governa ogni tipo
(\S\ref{cap:fiec}). Da qui in poi, nulla si introduce senza una di queste
cinque lavagne sotto.\upshape

\bigskip

% ---------------------------------------------------------------------
\chapter{Il punto di partenza}\label{cap:partenza}

{\spacedlowsmallcaps{Il formalismo fissa qualcosa che era gi\`a l\`i}}, non crea, cattura. Le regole non rendono vere le cose --- le cose sono vere indipendentemente dalle regole. Le regole sono uno strumento per essere espliciti su \emph{cosa si assume} e \emph{cosa si deriva}.

Questa distinzione \`e l'asse di tutto il libro. Ogni volta che il terreno smette di essere derivabile e diventa assunto, il testo lo segnala con un \emph{cero}: una proposizione numerata, scritta in stile sobrio, che enuncia l'assunzione e dice perch\'e regge anche senza dimostrazione interna. I ceri sono raccolti in fondo (\S\ref{cap:ceri}) e si possono leggere in sequenza per avere il riassunto completo delle assunzioni di MLTT.

\bigskip

\noindent La metafora del cero votivo non \`e ornamento. Un cero si accende perch\'e qualcosa --- non importa cosa --- regga: l'atto di accenderlo non costituisce prova, e tuttavia non \`e capriccio, perch\'e ci\`o che si \`e affidato al cero ha un peso reale. Nel formalismo accade lo stesso. Certi giudizi --- l'esistenza del metalinguaggio, l'esistenza dei contesti, lo statuto delle regole di computazione, la decisione di lasciare entrare termini dentro i tipi --- non si possono dimostrare dall'interno del sistema. Si accendono come ceri: si dichiarano, si tengono in vista, e si va avanti.

\bigskip

\noindent\textbf{Cosa non sono i ceri.} Non sono assiomi nel senso di Hilbert (proposizioni del linguaggio scelte come punto di partenza per derivare altre proposizioni del linguaggio): vivono nel metalinguaggio, o sono scelte di design. Non sono congetture (non si tratta di affermazioni che potrebbero un giorno rivelarsi false). Non sono postulati arbitrari: sono i punti in cui la materia del sistema chiede di essere fissata, e ogni fissaggio \`e una scelta motivata ma non costretta.

\bigskip

\noindent\textbf{Quanti sono.} Quattro nel corpo del libro, pi\`u un quinto al limite (la coerenza del sistema, indimostrabile dall'interno per Gödel). Sono pochi perch\'e MLTT \`e parsimoniosa. Una volta accesi, le quattro regole F-I-E-C (\S\ref{cap:fiec}) generano l'intero edificio.

% ---------------------------------------------------------------------
\chapter{I quattro giudizi del metalinguaggio}\label{cap:giudizi}

Tutto vive qui, fuori dal linguaggio oggetto. I quattro giudizi sono:

\begin{equation*}
\begin{array}{ll}
A : \Type & \text{\meta{$A$ \`e un tipo}} \\
a : A & \text{\meta{$a$ abita $A$}} \\
A \definitorio B : \Type & \text{\meta{$A$ e $B$ sono lo stesso tipo (uguaglianza definitoria)}} \\
a \definitorio b : A & \text{\meta{$a$ e $b$ sono lo stesso termine in $A$ (uguaglianza definitoria)}}
\end{array}
\end{equation*}

Questi \emph{non sono} tipi o termini: sono \textbf{affermazioni sul sistema}. Vivono nel metalinguaggio. Prima ancora di scrivere una regola, stai gi\`a assumendo che questo livello esista e abbia senso.

\begin{cero}[Esiste un metalinguaggio con i quattro giudizi]
Non lo si dimostra dall'interno del sistema: \`e la lavagna su cui si scrivono le regole. Prima assunzione, quella che rende possibili tutte le altre.
\end{cero}

\begin{cero}[Esistono i contesti e l'operazione $\vdash$]\label{cero:contesti}
``Assumere $x : A$'' \`e un'operazione legittima: i contesti si possono estendere, si pu\`o ragionare \emph{sotto} un'ipotesi. Senza, niente regole con premesse ipotetiche, niente tipi dipendenti.
\end{cero}

\begin{oss}[Sulle due uguaglianze]
L'uguaglianza $\definitorio$ di questa sezione \`e \textbf{definitoria}: un'asserzione sul sistema, non un tipo, non si dimostra. \`E quella che dice ``queste due cose riducono alla stessa cosa per le regole di computazione''. La incontreremo costantemente da qui a tutto il libro.

In \S\ref{cap:identita} introdurremo un \emph{altro} $\identita$, completamente diverso, che vive \textbf{nel linguaggio}: \`e un tipo, lo si abita con prove, lo si manipola come qualsiasi altro termine. Nei testi classici si scrivono uguali ($\equiv$ in entrambi i casi), e questa \`e una sciagura notazionale storica; in questo libro li distinguiamo tipograficamente apposta: $a \definitorio b : A$ per il primo, $a \identita b$ per il secondo, oppure $\mathsf{Id}\, A\, a\, b$ se serve massima chiarezza.
\end{oss}

% ---------------------------------------------------------------------
\chapter{Contesti e l'operazione $\vdash$}\label{cap:contesti}

\noindent {\spacedlowsmallcaps{Tutto quello che scriveremo da qui in poi vale \emph{sotto ipotesi}}}. Non c'\`e regola in MLTT --- nemmeno la formazione di $\mathbb{B}$ --- che non si svolga dentro un contesto. Era il Cero~\ref{cero:contesti}; resta da specificarne la grammatica.

\bigskip

\noindent\textbf{Il caso minimo dove serve.} Vogliamo poter scrivere la regola di formazione di $\Pi$ (\S\ref{cap:pi}):
\[
\inferrule{A : \Type \\ x : A \vdash B\,x : \Type}{(x : A) \to B\,x : \Type}
\]
La seconda premessa dice ``in un contesto in cui $x : A$ \`e assunto, $B\,x$ \`e un tipo''. Senza una nozione formale di contesto e di operazione $\vdash$, quella riga \`e gergo. Con la grammatica esplicita, \`e una regola.

\section*{Grammatica dei contesti}

I contesti sono sequenze finite di ipotesi della forma $x : A$. Li indichiamo con $\Gamma, \Delta$, eventualmente con apici. Introduciamo un nuovo giudizio del metalinguaggio:
\[
\Gamma \;\mathsf{ctx} \qquad \text{\meta{$\Gamma$ \`e un contesto ben formato}}
\]

\noindent Le regole di formazione del giudizio $\mathsf{ctx}$ sono due, e generano tutti e soli i contesti legittimi.

\paragraph{Contesto vuoto.}
\[
\inferrule{ }{\cdot \;\mathsf{ctx}}
\]
Il contesto vuoto, sempre ben formato.

\paragraph{Estensione.}
\[
\inferrule{\Gamma \;\mathsf{ctx} \\ \Gamma \vdash A : \Type \\ x \notin \mathrm{dom}(\Gamma)}{\Gamma, x : A \;\mathsf{ctx}}
\]
Si pu\`o estendere $\Gamma$ con un'ipotesi $x : A$ \emph{a patto che} $A$ sia un tipo derivabile in $\Gamma$ (terza premessa) e $x$ sia un nome nuovo, non gi\`a presente (quarta premessa, condizione di freschezza).

\bigskip

\noindent\textbf{Osservazioni meccaniche.}
\begin{itemize}
\item Il contesto cresce \emph{a destra}: $\Gamma, x : A$ si legge ``$\Gamma$ esteso con $x : A$''. Le ipotesi pi\`u vecchie stanno a sinistra.
\item Le ipotesi pi\`u recenti possono dipendere dalle precedenti --- il tipo $B\,x$ in $\Gamma, x : A, y : B\,x$ usa la $x$ del passo prima. Senza questa possibilit\`a non si avrebbero tipi dipendenti, ed \`e proprio questa l'apertura formale del Cero~\ref{cero:dipendenza} che vedremo in \S\ref{cap:pi}.
\item Il giudizio $\mathsf{ctx}$ \`e del \textbf{metalinguaggio}, come gli altri quattro; vive sopra il linguaggio.
\end{itemize}

\section*{I quattro giudizi, riscritti sotto contesto}

I quattro giudizi di \S\ref{cap:giudizi} non compaiono mai \emph{nudi} nelle regole reali: portano sempre un contesto a sinistra di $\vdash$.

\begin{equation*}
\begin{array}{ll}
\Gamma \vdash A : \Type & \text{\meta{nel contesto $\Gamma$, $A$ \`e un tipo}} \\
\Gamma \vdash a : A & \text{\meta{nel contesto $\Gamma$, $a$ abita $A$}} \\
\Gamma \vdash A \definitorio B : \Type & \text{\meta{$A$ e $B$ sono lo stesso tipo in $\Gamma$}} \\
\Gamma \vdash a \definitorio b : A & \text{\meta{$a$ e $b$ sono lo stesso termine di $A$ in $\Gamma$}}
\end{array}
\end{equation*}

\noindent Il contesto non \`e ``parte'' del giudizio: \`e \emph{l'ambiente di ipotesi sotto cui il giudizio vale}. Contesti diversi producono giudizi diversi della stessa forma: $\cdot \vdash a : A$ (sotto vuoto) e $x : B \vdash a : A$ sono due asserzioni distinte, anche con $a$ e $A$ identici.

\section*{La regola della variabile}

Senza una regola che permetta di \emph{usare} le ipotesi, il contesto sarebbe decorativo. La regola che le mette in gioco \`e la pi\`u semplice possibile:

\[
\inferrule{\Gamma, x : A, \Delta \;\mathsf{ctx}}{\Gamma, x : A, \Delta \vdash x : A}
\]

\noindent Letta: ``se nel contesto compare $x : A$, allora posso derivare $x : A$''. \`E ovvia; \`e quella che giustifica formalmente l'idea ``ho $x$ a disposizione perch\'e l'ho assunto''. \`E anche l'unica regola che ``guarda dentro'' il contesto --- tutte le altre lo lasciano com'\`e o lo estendono in modo prevedibile.

\section*{Le altre regole strutturali}

Tre regole che governano \emph{come} il contesto pu\`o cambiare lungo una derivazione senza alterarne il significato. Le scriviamo nella loro forma minima; lo scope della trattazione le tiene a questo livello: \emph{The Little Typer} e \emph{Verified Functional Programming in Agda} non chiedono di pi\`u. Nessuna di esse \`e un cero: sono \emph{lemmi ammissibili}, conseguenze del modo in cui le regole specifiche dei tipi (Parte~II) saranno scritte. Si dichiarano qui perch\'e da Parte~II in poi le useremo tacitamente.

\paragraph{Indebolimento (\emph{weakening}).}
\[
\inferrule{\Gamma, \Delta \vdash \mathcal{J} \\ \Gamma \vdash A : \Type \\ x \notin \mathrm{dom}(\Gamma, \Delta)}{\Gamma, x : A, \Delta \vdash \mathcal{J}}
\]
Aggiungere un'ipotesi non rovina ci\`o che era gi\`a derivato: se sapevi $\mathcal{J}$, lo sai ancora con un'ipotesi in pi\`u. Qui $\mathcal{J}$ sta per uno qualsiasi dei quattro giudizi --- usiamo la calligrafica $\mathcal{J}$ per non confonderla con l'eliminatore $\Jelim$ del tipo identit\`a (\S\ref{cap:identita}).

\paragraph{Sostituzione.}
\[
\inferrule{\Gamma \vdash a : A \\ \Gamma, x : A, \Delta \vdash \mathcal{J}}{\Gamma, \Delta[a/x] \vdash \mathcal{J}[a/x]}
\]
Se in $\Gamma$ esibisci un abitante $a$ di $A$, ogni derivazione fatta \emph{sotto l'ipotesi $x : A$} si pu\`o specializzare a quel concreto $a$, sostituendolo dappertutto --- nelle ipotesi successive che dipendevano da $x$ e nel giudizio finale. \`E la regola che chiude l'ipotesi.

\paragraph{Conversione di tipo.}
\[
\inferrule{\Gamma \vdash a : A \\ \Gamma \vdash A \definitorio B : \Type}{\Gamma \vdash a : B}
\]
Se $A$ e $B$ sono definitoriamente uguali, abitare $A$ e abitare $B$ \`e la stessa cosa. \`E la regola che fa lavorare $\definitorio$ \emph{dentro} il sistema dei tipi: senza, l'uguaglianza definitoria si fermerebbe a essere un'asserzione esterna e non avrebbe effetto operativo sul controllo dei tipi.

\section*{Convenzione tipografica per il resto del libro}

Da \S\ref{cap:fiec} in poi, le regole di inferenza si scrivono \emph{senza contesto in vista} quando tutte le premesse e la conclusione condividono lo stesso. Si sottintende che ce ne sia uno, lo stesso ovunque nella regola. Esempio: la regola di introduzione di $\mathbb{B}$ (\S\ref{cap:bool})
\[
\inferrule{ }{\tto : \mathbb{B}}
\]
sta in realt\`a per
\[
\inferrule{\Gamma \;\mathsf{ctx}}{\Gamma \vdash \tto : \mathbb{B}}
\]
per qualunque $\Gamma$ ben formato.

\bigskip

\noindent Quando invece una premessa \emph{estende} il contesto rispetto alla conclusione --- caso paradigmatico, la $\Pi$-formazione di \S\ref{cap:pi}:
\[
\inferrule{A : \Type \\ x : A \vdash B\,x : \Type}{(x : A) \to B\,x : \Type}
\]
--- il $\vdash$ sopravvive nella scrittura, perch\'e \`e l\`i che il contesto sta facendo un lavoro non banale. La regola, sotto contesto generico $\Gamma$, sarebbe per esteso:
\[
\inferrule{\Gamma \vdash A : \Type \\ \Gamma, x : A \vdash B\,x : \Type}{\Gamma \vdash (x : A) \to B\,x : \Type}
\]
La compressione ``ometti $\Gamma$ dove \`e ovvio, lascialo dove non lo \`e'' \`e un fatto di leggibilit\`a; sotto, la regola \`e la seconda.

\section*{Test diretto}

\begin{oss}[$\vdash$ e ``,'' sono entrambi del metalinguaggio]
Il ``,'' che separa le ipotesi in $\Gamma, x : A$ \`e grammatica dei contesti, non del linguaggio: non puoi passare $\Gamma, x : A$ come argomento di una funzione, perch\'e non \`e un oggetto del linguaggio. \`E un contesto, e i contesti vivono nel metalinguaggio. Stessa cosa per $\vdash$ --- gi\`a sotto test in \S\ref{cap:test}, ma ora con grammatica precisa alle spalle.
\end{oss}

\bigskip

\noindent Con la grammatica di $\vdash$ in tasca, ogni regola del seguito ha un significato preciso e non interpretativo. Ogni volta che leggi una premessa con un nuovo $x : A$, sai cosa si sta facendo: estensione di contesto via la regola di estensione, e da quel momento $x$ \`e disponibile come variabile via la regola della variabile.

% ---------------------------------------------------------------------
\chapter{Linguaggio e metalinguaggio --- il test}\label{cap:test}

Sezione operativa: ogni volta che sui livelli ci si perde, si torna qui.

\section*{Il test in una riga}

\noindent\itshape Questo oggetto lo posso passare come argomento a una funzione, o ci posso costruire sopra un altro termine?\upshape

\medskip

\noindent \textbf{S\`i} $\Rightarrow$ vive nel \textbf{linguaggio}. \`E un termine o un tipo: oggetto manipolabile.

\noindent \textbf{No} --- \`e qualcosa che \emph{autorizza} a scrivere un giudizio ma non \`e esso stesso un oggetto manipolabile --- $\Rightarrow$ vive nel \textbf{metalinguaggio}. \`E una regola, un giudizio, un'affermazione \emph{sul} sistema.

\bigskip

\noindent In altre parole: il linguaggio \`e ci\`o di cui parli; il metalinguaggio \`e ci\`o \emph{con cui} parli. \texttt{zero} \`e una cosa di cui parli (linguaggio); ``\texttt{zero} $:\mathbb{N}$'' \`e una cosa che \emph{dici} --- un'affermazione (metalinguaggio).

\section*{Tre punti dove si inciampa}

\paragraph{(i) I due punti \texttt{:} non sono dentro il linguaggio.}
\texttt{a : A} non \`e un oggetto unico. \texttt{a} e \texttt{A} vivono nel linguaggio; il \texttt{:} che li unisce \`e un'\emph{affermazione sul sistema} (``asserisco che \texttt{a} abita \texttt{A}''). Non si pu\`o passare \texttt{a : A} come argomento a niente: non \`e un oggetto, \`e un giudizio. Il \texttt{:} \`e metalinguaggio puro.

\paragraph{(ii) La barra orizzontale \rule{1cm}{0.4pt} \`e metalinguaggio.}
Separa premesse da conclusione: ``dati questi giudizi sopra, ottieni questo giudizio sotto''. \`E la grammatica con cui si scrivono le regole (Cero~1).

\paragraph{(iii) Il $\vdash$ \`e metalinguaggio, ma fa una cosa diversa dalla barra.}
Sta \emph{dentro una singola premessa}: $x : A \vdash B\,x : \Type$, ``nel contesto in cui assumo $x : A$, vale $B\,x : \Type$''. Costruisce \emph{un} giudizio dicendo sotto quali ipotesi vale (Cero~\ref{cero:contesti}; la grammatica completa di contesti e $\vdash$ \`e in \S\ref{cap:contesti}). La barra combina giudizi interi; il $\vdash$ costruisce un giudizio sotto ipotesi. Non sono lo stesso simbolo.

\bigskip

\noindent Il quarto punto dove si inciampa --- la coppia $\definitorio$ vs.\ $\identita$ --- non lo si pu\`o ancora introdurre per filo perch\'e ci manca il tipo identit\`a; lo chiuderemo in \S\ref{cap:identita}. Per ora basti la promemoria: il $\definitorio$ visto in \S\ref{cap:giudizi} \`e definitorio, vive nel metalinguaggio, non si abita; il $\identita$ che incontreremo in \S\ref{cap:identita} \`e un tipo del linguaggio, si abita con \texttt{refl}.

\section*{Tabella lampo}

\begin{center}
\renewcommand{\arraystretch}{1.15}
\begin{tabular}{lll}
\toprule
\textit{Vedi questo} & \textit{Dov'\`e} & \textit{Perch\'e} \\
\midrule
\texttt{tt}, \texttt{ff}, \texttt{zero}, \texttt{suc n} & linguaggio & termini, manipolabili \\
$\mathbb{B}$, $\mathbb{N}$, \texttt{Vec A n} & linguaggio & tipi, oggetti del sistema \\
\texttt{:} in \texttt{a : A} & meta & giudizio, non oggetto \\
$\definitorio$ definitorio (\S\ref{cap:giudizi}) & meta & asserzione sul sistema \\
\rule{0.9cm}{0.3pt} (la barra) & meta & combina giudizi interi (Cero~1) \\
$\vdash$ & meta & giudizio sotto ipotesi (Cero~\ref{cero:contesti}) \\
coerenza del sistema & meta-meta & non dall'interno --- G\"odel \\
\bottomrule
\end{tabular}
\end{center}

\noindent Le righe sui tipi del linguaggio le completiamo in \S\ref{cap:identita}, quando il tipo identit\`a porter\`a con s\'e altri abitanti del linguaggio.

% ---------------------------------------------------------------------
\chapter{Lo schema F-I-E-C}\label{cap:fiec}

\noindent {\spacedlowsmallcaps{Questa \`e la sezione cardine}}. Ogni tipo in MLTT \`e dato da \textbf{quattro regole}, sempre nello stesso ordine, con la stessa struttura. Non c'\`e altro modo di introdurre un tipo, e non c'\`e altro modo di lavorarci sopra. Tutto quello che si vedr\`a nella Parte~II \`e declinazione di questo schema.

\bigskip

\begin{center}
\renewcommand{\arraystretch}{1.25}
\begin{tabular}{lll}
\toprule
\textit{Regola} & \textit{Cosa dice} & \textit{Dove vive} \\
\midrule
\textbf{Formazione} (F) & quando il tipo esiste & linguaggio \\
\textbf{Introduzione} (I) & come costruisco un abitante (costruttori) & linguaggio \\
\textbf{Eliminazione} (E) & come uso un abitante (eliminatori) & \textbf{metalinguaggio} \\
\textbf{Computazione} (C) & come riduce (regole $\beta$, definitorie) & \textbf{metalinguaggio} \\
\bottomrule
\end{tabular}
\end{center}

\bigskip

Due cose vanno fissate subito, perch\'e ricorrono identiche per ogni tipo.

\section*{Gli eliminatori non sono funzioni}

Non sono funzioni che abitano un tipo: sono \emph{regole} del sistema di deduzione. Quando un testo li presenta come oggetti applicabili ($\indN(\dots)$, $\indB(\dots)$) sta facendo un'\textbf{immersione pedagogica} --- porta qualcosa dal metalinguaggio dentro il linguaggio. L'oggetto cambia natura nel farlo. Da tenere a mente ogni volta che si incontra un eliminatore scritto come se fosse un termine.

\begin{oss}[Test diretto]
\agdainl{suc} si pu\`o passare a una funzione higher-order? S\`i, dunque \`e nel linguaggio. \agdainl{ind-N} in MLTT pura si pu\`o passare? No, \`e una regola del metalinguaggio. Agda --- per ergonomia --- ti far\`a scrivere una funzione di nome \texttt{nat-rec} che si comporta come \agdainl{ind-N}; quella \`e l'immersione pedagogica resa cittadino del linguaggio. Non confondere il gesto pedagogico con la regola.
\end{oss}

\section*{Le regole di computazione sono definitorie}

Non sono teoremi. Non si ha nulla per dimostrarle \emph{dall'interno} del sistema. Si scrivono perch\'e rispecchiano la struttura del tipo. Il sistema funziona perch\'e \`e allineato con quella struttura, non per decreto.

\begin{cero}[Le regole di computazione sono definitorie]\label{cero:comp}
Si assumono, e reggono perch\'e la struttura era gi\`a l\`i.
\end{cero}

\section*{Forma generale di una regola}

Tutte le regole hanno la forma:
\[
\inferrule{\text{premessa}_1 \\ \text{premessa}_2 \\ \dots}{\text{conclusione}}
\]

Premesse e conclusione sono \emph{giudizi} del metalinguaggio (\texttt{a : A}, \texttt{A : Type}, $x : A \vdash B\,x : \Type$, ecc.). La barra dice ``se hai i giudizi sopra, puoi scrivere il giudizio sotto''. Una sequenza di regole applicate fino a un giudizio senza premesse aperte \`e una \emph{derivazione}.

\bigskip

\noindent Da \S\ref{cap:bool} in poi, ogni tipo riceve esattamente queste quattro regole, scritte in questo formato, in quest'ordine.

\section*{Stile di esposizione: alla Little Typer}\label{sec:stile}

\noindent {\spacedlowsmallcaps{Una dichiarazione di metodo}}, prima di entrare in Parte~II.

\bigskip

Esistono due modi storicamente attestati di presentare un tipo induttivo. Il primo --- quello di \emph{The Little Typer} di Friedman e Christiansen, e degli articoli fondativi di Martin-L\"of --- elenca per esteso F-I-E-C e usa \textbf{soltanto} l'eliminatore esplicito ($\indB$, $\indN$, $\Jelim$, $\ldots$) quando arriva il momento di costruire abitanti. Il secondo --- quello di Agda --- scrive una \texttt{data} \texttt{where} per dichiarare il tipo, e poi tutto il resto via \emph{pattern matching}: clausole con costruttori a sinistra del segno uguale, niente eliminatore in vista.

I due stili \emph{non sono equivalenti}: il secondo \`e pi\`u potente, e quanto pi\`u potente lo discuteremo in apertura di Parte~III (\S\ref{cap:prima}). Per ora ci basta la scelta di forma:

\begin{itemize}
\item \textbf{In questo libro le regole le scriviamo alla Little Typer.} Ogni tipo riceve F, I, E, C in forma di regola di inferenza esplicita; l'eliminatore ha un nome ($\indB$, $\indN$, $\ldots$) e si presenta come regola del metalinguaggio.
\item \textbf{Quando arrivano i sorgenti Agda} (Parte~III), useremo il pattern matching come si presenta nei libri di riferimento --- a cominciare da \emph{Verified Functional Programming in Agda} di Stump. Il nostro lavoro di lettura sar\`a sempre: \emph{quale F-I-E-C \`e in atto sotto questa clausola}?
\end{itemize}

\noindent\textbf{Esercizio centrale.} Una volta visti i sei tipi di Parte~II, il lettore dovrebbe essere in grado di guardare una funzione Agda definita per casi e ricostruirne meccanicamente la versione con eliminatore esplicito. La capacit\`a di fare quest'operazione --- almeno nei casi facili --- \`e il segnale che l'F-I-E-C \`e diventata vista, non solo lessico.